% ID: SEL-D0004-OA-E01
% ESTADO: cotejada. Cotejo literal independiente y revisión de composición realizados.
% Fuente: Selectividad/sel_2019_matematicas_aplicadas/Suplente Septiembre/Exámenes Matemáticas Aplicadas CCSS - A y B.pdf
% SHA-256: eecad753694b3d8cb2d427efddb00d75faae9c4ec5b027056291bc95d7d2e3da
% Página: 1; opción A; ejercicio 1.
% Encabezado original: EJERCICIO 1. Solo enunciado; no contiene solución.

Una granja elabora una dieta mezclando dos tipos de pienso A y B. El pienso
A aporta 2 unidades de Calcio y 1 de Hierro por cada kilogramo, mientras
que el B aporta 1 de Calcio y 2 de Hierro. El coste por kilogramo tanto del
pienso A como del pienso B es 1 euro por kilogramo. La dieta deberá aportar
al menos 2 unidades de Calcio y 2 de Hierro.

Determine los kilogramos que se han de mezclar de cada tipo de pienso para
que el coste de la dieta sea mínimo. \textquestiondown{}Cuál sería dicho
coste? \textquestiondown{}Cuántas unidades de Hierro y de Calcio se
administrarían a los animales con esta dieta?
